\cleardoublepage

\section{绪论}

\subsection{研究背景}

近年来，大语言模型（Large Language Models, LLM）在问答、摘要、对话等任务中表现出强大的生成能力，但在知识截止、事实性与可控性方面仍存在不足，尤其是在需要严格依据外部证据的场景中容易产生“幻觉”\cite{gpt3, gpt4}。检索增强生成（Retrieval-Augmented Generation, RAG）通过将检索模块与生成模块结合，使模型在回答前能够从外部知识库中检索证据，从而提升事实性与可解释性\cite{ref1}。

然而，在真实使用中，RAG系统面临更复杂的输入分布：用户查询可能包含语义歧义（如“Python”既指编程语言也指蛇）、否定或约束条件（如“not named python”），也可能出现知识库中根本不存在答案的情况。若检索与重排序策略不能鲁棒处理上述场景，系统会返回不相关证据并诱发生成端的错误回答，从而影响用户体验与系统可信度\cite{ref26}。

\subsection{研究问题}

本文的研究目标并不是抽象地讨论“更强的RAG”，而是围绕一个可复现的系统原型回答三个更具体的问题。首先，在本科论文可承受的算力和数据规模内，如何搭建一套端到端实验平台，使检索、重排序、补检索与验证模块都能在统一入口下比较，而不是分散在难以复核的脚本里。其次，当查询包含歧义词、约束条件或分散在不同段落中的多跳线索时，简单的一轮召回往往无法覆盖完整证据，那么重排序、结果融合和自适应补检索究竟能在多大程度上改善Top-$k$覆盖率。最后，当检索到的证据本身不完整甚至不存在时，系统是否应该拒答，以及拒答机制应当怎样设置，才能在抑制幻觉和维持可用性之间取得相对稳妥的平衡。现有工作往往分别优化检索、生成或验证模块，但缺乏对多模块组合下系统级行为的系统性分析，尤其缺少对失败模式的定量刻画。以上三个问题共同指向题目中的“多源异构知识融合”和“复杂推理机制”：问题一提供统一证据入口与实验载体，问题二考察多源证据如何影响复杂多跳检索覆盖，问题三则进一步分析生成与验证模块如何改变最终可用性。

\subsection{本文工作}

围绕前述问题，本文将系统实现作为机制研究的可复现实验载体，而不是单纯的工程集成。系统层面的第一个工作，是把文本段落、半结构化表格以及图谱三元组统一转写为可检索的证据单元，并放在同一套检索与重排序流程中处理，以便观察异构知识是否真的能够补足多跳问题中的信息缺口\cite{tapas, tabert}。这一设计在论文中被抽象为统一证据接口（Unified Evidence Interface）：它把“证据来自哪种载体”和“证据如何进入检索流水线”解耦，使后续模块可以在同一接口上比较文本、表格和图谱证据。第二个工作，是设计基于重排序分数的补检索门控：当首轮候选的置信度偏低时，系统再触发一次扩展检索；当首轮结果已经足够明确时，就避免额外的查询开销。本文将该机制定位为面向召回率的检索控制器，而不是直接优化最终答案分数的生成策略\cite{ref11}。

第三个工作，是对召回证据做轻量整理，把原本分散的文本、表格行和图谱事实按实体与来源顺序拼接成更容易被生成模型利用的上下文。这里所说的“复杂推理”，并不是系统内部真的实现了一套显式符号推理器，而是指系统尽量把多跳问题需要的材料准备完整，让生成模型能够在较好的上下文基础上完成隐式推理。第四个工作，是在生成之后增加事实核验与拒答环节，对回答中的声明逐条做证据支持检查\cite{ref24}。如果证据明显缺失或出现冲突，系统就不再继续输出答案，而是转为拒答。

为了判断这些设计的真实作用，本文没有停留在模块说明层面，而是安排了四组相互衔接的实验：先用全量消融看各模块的净影响，再用更强基线与跨数据集比较观察Adaptive策略的相对位置，随后通过阈值扫描讨论部署代价，最后结合错误分布和案例分析解释系统为什么会成功、又为什么会失败。进一步地，本文把验证模块在严格配置下出现的大规模拒答现象概括为 verification collapse，即验证器在噪声证据和硬阈值共同作用下将大量可回答样本压成拒答。围绕这一现象，后续实验还设计了 Verification-Aware Retrieval (VAR)，用于检验“验证器只拒答”和“验证器回写检索”之间的差异。换言之，本文从系统组合角度研究多模块 RAG 的耦合效应，并重点分析其失败模式与优化路径，而不是仅仅给出一组更高的数字。

本文的主要贡献总结如下：
\begin{enumerate}
    \item 提出统一证据接口（Unified Evidence Interface），解耦异构模态与检索流程，并揭示格式噪声问题；
    \item 证明自适应检索仅提升召回，无法线性传导至端到端生成；
    \item 首次发现并定义验证崩溃（Verification Collapse），指出其本质来源于“噪声证据 $\times$ 硬阈值决策”的误差放大效应，并提出 VAR 闭环修复。
\end{enumerate}

\subsection{相关工作 (Related Work)}

现有RAG研究大致可以分成四条与本文直接相关的路线。第一类工作聚焦纯文本检索与重排序。以DPR\cite{ref6}为代表的系统已经证明，稠密检索配合交叉编码器重排序可以显著改善开放域问答中的证据质量\cite{bm25, bert_rerank}。不过这类方法主要面向非结构化文本，面对表格字段或图谱关系时往往需要额外的表示转换\cite{manning_ir}。

第二类工作尝试把结构化知识纳入RAG。HybridQA\cite{hybridqa}、OTT-QA\cite{ottqa}等研究已经说明，文本和表格往往需要联合使用；Think-on-Graph\cite{ref17}、HippoRAG\cite{ref14}等工作则进一步把图谱关系引入检索过程\cite{graphdatabases, ref29}。这些研究为本文提供了很重要的启发，但它们也暴露出一个现实问题：一旦把多种载体放进同一套流程，表示不统一、排序信号相互干扰的工程成本会迅速上升\cite{ref30}。

第三类工作关注主动式检索，也就是让系统在生成过程中决定“什么时候还要再查一次”。这类方法对HotpotQA\cite{ref22}等需要多跳证据的问题尤其重要，因为单次检索常常只能找到第一跳线索\cite{twiki}。FLARE\cite{ref12}、IRCoT\cite{ref26}等代表性方法展示了主动检索的潜力\cite{ref10, ref23}，但它们通常伴随更高的Token开销和推理延迟\cite{ref33, ref34}。本文的Adaptive设计实际上也属于这一路线，只是采用了更轻量、也更容易复核的门控方式。

第四类工作围绕生成后的自我校验和防幻觉展开。Self-RAG\cite{ref8}、CoVe\cite{ref24}等方法都试图让系统在输出之前再做一次检查，从而降低无依据生成的概率\cite{ref25, ref35}。但从工程角度看，验证器并不天然等于更可靠的系统；如果判定逻辑过于刚性，验证步骤本身也可能变成新的误差来源。本文后续实验里最重要的发现之一，正是这一点。

因此，本文的研究空白在于：在统一、可复现的实验载体中，把“异构知识接入”“主动补检索”“生成后验证”三件常被分别讨论的机制放到同一流程里比较，并据此分析它们之间的真实耦合关系与误差传导路径。

\subsection{论文结构}

本文后续结构安排如下：第2章介绍系统总体设计与关键算法；第3章给出实验设计与评估指标；第4章展示实验结果与分析；最后总结全文并讨论未来工作方向。
